\documentclass[../DoAn.tex]{subfiles}
\begin{document}

\section*{Quy định chung}

\section{Ngành học}

\section{Đánh dấu (bullet) và đánh số (numering)}

\section{Cách thêm bảng}

\section{Chèn hình ảnh}

\section{Tài liệu tham khảo}

\paragraph{Cách liệt kê}\mbox{}

\paragraph{Các loại tài liệu tham khảo}\mbox{}

\section{Cách viết phương trình và công thức toán học}

\section{Qui cách đóng quyển}

% Nguồn: D:/hust-master/Thesis/SUMO/rules/bounds/param_bounds.yaml và rules/specs/B1_manual.json.
% Giá trị chép nguyên từ tệp, chỉ đổi dấu thập phân sang dấu phẩy.
\section{Miền giá trị tham số vType}
\label{section:appendix_bounds}

Bảng~\ref{table:appendix_bounds_1} và Bảng~\ref{table:appendix_bounds_2} liệt kê toàn bộ 24 thuộc tính vType trong tệp miền giá trị dùng chung cho LLM và Optuna (Mục~\ref{section:3.3}). Cột ``Mô hình'' ghi mô hình con của SUMO mà thuộc tính thuộc về: \texttt{vType} (thuộc tính chung của kiểu xe), \texttt{carFollow} (mô hình bám xe Krauss), \texttt{LC} (mô hình chuyển làn SL2015) và \texttt{junction} (mô hình nút giao). Cột ``Mặc định'' là giá trị mà vType của B0 mang khi đặc tả không khai báo thuộc tính, tức giá trị mặc định của SUMO theo vClass; giá trị này có thể nằm ngoài $[\min,\max]$ khi nó là giá trị đặc biệt mang nghĩa ``tắt'' và bộ kiểm tra hợp lệ không kiểm tra nó. Các khoảng áp dụng chung cho ba vClass cơ sở (\texttt{passenger}, \texttt{motorcycle}, \texttt{custom1}) trừ ba thuộc tính có khoảng riêng theo vClass: với \texttt{motorcycle} và \texttt{custom1}, \texttt{minGap} có miền $[0{,}5;\,2{,}5]$~m, mặc định 2,5~m; \texttt{accel} có miền $[0{,}8;\,6{,}0]$~m/s\textsuperscript{2}, mặc định 6,0; \texttt{decel} có miền $[2{,}0;\,10{,}0]$~m/s\textsuperscript{2}, mặc định 10,0. Các dòng có dấu (*) trong bảng là các thuộc tính này; giá trị ghi trong bảng là của \texttt{passenger}.

\begin{table}[H]
\centering
\small\setlength{\tabcolsep}{3pt}
\begin{tabular}{|p{3.0cm}|p{1.6cm}|p{1.2cm}|p{1.1cm}|p{1.1cm}|p{1.6cm}|p{3.2cm}|}
\hline
\textbf{Tham số} & \textbf{Mô hình} & \textbf{Đơn vị} & \textbf{Min} & \textbf{Max} & \textbf{Mặc định} & \textbf{Nhóm hành vi} \\ \hline
\texttt{tau} & carFollow & s & 0,4 & 3,0 & 1,0 & \texttt{f\_\allowbreak risk} \\ \hline
\texttt{minGap} (*) & vType & m & 1,0 & 3,0 & 2,5 & \texttt{encroachment}, \texttt{illegal\_\allowbreak left\_\allowbreak box} \\ \hline
\texttt{accel} (*) & carFollow & m/s\textsuperscript{2} & 0,8 & 4,0 & 2,6 & \texttt{I\_\allowbreak slow} \\ \hline
\texttt{decel} (*) & carFollow & m/s\textsuperscript{2} & 2,0 & 7,5 & 4,5 & \texttt{f\_\allowbreak risk} \\ \hline
\texttt{sigma} & carFollow & -- & 0,0 & 1,0 & 0,5 & \texttt{I\_\allowbreak slow} \\ \hline
\texttt{speedFactor} & vType & -- & 0,5 & 1,3 & 1,0 & \texttt{I\_\allowbreak slow} \\ \hline
\texttt{speedDev} & vType & -- & 0,0 & 0,3 & 0,1 & \texttt{I\_\allowbreak slow} \\ \hline
\texttt{minGapLat} & LC & m & 0,3 & 1,0 & 0,6 & \texttt{encroachment} \\ \hline
\texttt{maxSpeedLat} & vType & m/s & 0,5 & 3,0 & 1,0 & \texttt{encroachment} \\ \hline
\texttt{latAlignment} & vType & -- & \multicolumn{2}{p{2.2cm}|}{giá trị rời rạc: \texttt{right}, \texttt{center}, \texttt{left}, \texttt{compact}, \texttt{nice}, \texttt{arbitrary}} & \texttt{center} & \texttt{encroachment} \\ \hline
\texttt{lcSublane} & LC & -- & 0,0 & 5,0 & 1,0 & \texttt{encroachment} \\ \hline
\texttt{lcPushy} & LC & -- & 0,0 & 1,0 & 0,0 & \texttt{encroachment} \\ \hline
\end{tabular}
\caption{Miền giá trị tham số vType (phần 1: bám xe, thuộc tính chung và làn con)}
\label{table:appendix_bounds_1}
\end{table}

\begin{table}[H]
\centering
\small\setlength{\tabcolsep}{3pt}
\begin{tabular}{|p{3.0cm}|p{1.6cm}|p{1.2cm}|p{1.1cm}|p{1.1cm}|p{1.6cm}|p{3.2cm}|}
\hline
\textbf{Tham số} & \textbf{Mô hình} & \textbf{Đơn vị} & \textbf{Min} & \textbf{Max} & \textbf{Mặc định} & \textbf{Nhóm hành vi} \\ \hline
\texttt{lcPushyGap} & LC & m & 0,0 & 1,0 & 0,6 & \texttt{f\_\allowbreak risk} \\ \hline
\texttt{lcAssertive} & LC & -- & 0,5 & 5,0 & 1,0 & \texttt{f\_\allowbreak risk} \\ \hline
\texttt{lcImpatience} & LC & -- & 0,0 & 1,0 & 0,0 & \texttt{patience} \\ \hline
\texttt{lcTimeTo\allowbreak Impatience} & LC & s & 1,0 & 600,0 & $\infty$ & \texttt{patience} \\ \hline
\texttt{lcAccelLat} (cố định) & LC & m/s\textsuperscript{2} & 1,0 & 1,0 & 1,0 & \texttt{encroachment} \\ \hline
\texttt{lcMaxSpeed\allowbreak LatStanding} & LC & m/s & 0,0 & 3,0 & 1,0 & \texttt{encroachment} \\ \hline
\texttt{lcSpeedGain} & LC & -- & 0,0 & 10,0 & 1,0 & \texttt{encroachment} \\ \hline
\texttt{impatience} & junction & -- & 0,0 & 1,0 & 0,0 & \texttt{patience} \\ \hline
\texttt{jmDriveAfter\allowbreak RedTime} & junction & s & 0,0 & 10,0 & $-1{,}0$ & \texttt{patience} \\ \hline
\texttt{jmDriveRed\allowbreak Speed} & junction & m/s & 1,0 & 13,9 & 55,56 (**) & \texttt{patience} \\ \hline
\texttt{jmDriveAfter\allowbreak YellowTime} (cố định) & junction & s & $-1{,}0$ & $-1{,}0$ & $-1{,}0$ & \texttt{patience} \\ \hline
\texttt{jmTimegap\allowbreak Minor} (chỉ \texttt{custom1}) & junction & s & 0,5 & 3,0 & 1,0 & \texttt{illegal\_\allowbreak left\_\allowbreak box}, \texttt{f\_\allowbreak risk} \\ \hline
\end{tabular}
\caption{Miền giá trị tham số vType (phần 2: chuyển làn SL2015 và mô hình nút giao)}
\label{table:appendix_bounds_2}
\end{table}

\begin{sloppypar}
Ghi chú cho Bảng~\ref{table:appendix_bounds_2}: (**) giá trị mặc định của \texttt{jmDriveRedSpeed} trong tệp là 55,55555555555556~m/s, bằng tốc độ tối đa \texttt{maxSpeed} (200~km/h) của \texttt{passenger} và \texttt{motorcycle}. Mặc định $\infty$ của \texttt{lcTimeToImpatience} tắt cơ chế tăng dần mức thiếu kiên nhẫn; mặc định $-1$ của \texttt{jmDriveAfterRedTime} nghĩa là không bao giờ đi khi đèn đỏ, và muốn giữ các mặc định này thì đặc tả phải bỏ trống thuộc tính. Mặc định của \texttt{lcPushyGap} bằng \texttt{minGapLat} và của \texttt{lcMaxSpeedLatStanding} bằng \texttt{maxSpeedLat}. Hai thuộc tính \texttt{lcAccelLat} và \texttt{jmDriveAfterYellowTime} được cố định ($\min = \max =$ mặc định) vì không có chỉ số nào ràng buộc chúng. \texttt{jmTimegapMinor} chỉ được đặt cho \texttt{custom1} vì chỉ chỉ số M7 ràng buộc nó.
\end{sloppypar}

Hai thuộc tính tồn tại trong SUMO 1.27 nhưng bị loại khỏi không gian tham số là \texttt{jmIgnoreFoeProb} và \texttt{jmIgnoreFoeSpeed}. Thuộc tính thứ nhất khiến xe bỏ qua xe xung đột có quyền ưu tiên và do đó tạo va chạm thật theo thiết kế; thuộc tính thứ hai chỉ có tác dụng cùng thuộc tính thứ nhất. Bộ kiểm tra hợp lệ báo hai thuộc tính này với mã lỗi \texttt{unknown\_attr} kèm lý do loại.

\section{Đặc tả luật B1}
\label{section:appendix_b1}

Bộ luật thủ công B1 (\texttt{spec\_id} \texttt{B1\_manual}, nguồn \texttt{manual}) được chuyển giao từ mô hình giao thông hỗn hợp ô tô -- xe máy của Saval và cộng sự~\cite{saval2023mixed}, được hiệu chỉnh trên video đường Tây Sơn, Hà Nội. Các tham số bám xe lấy từ Bảng 11 của bài báo (tham số tài xế cho Hà Nội, mô hình IDM) và được ánh xạ sang mô hình Krauss của SUMO theo cùng ý nghĩa: gia tốc cực đại $a$ thành \texttt{accel}, giảm tốc cực đại $b$ thành \texttt{decel}, khoảng thời gian an toàn $T$ thành \texttt{tau} và khoảng cách an toàn khi dừng $s_0$ thành \texttt{minGap}; do khác mô hình, các giá trị này mang nhãn chuyển giao. Cấu trúc của B1 (ba nguyên mẫu, các tham số được khai báo và hai luật) là cấu trúc cố định mà trong đó các phương án tối ưu dự kiến điều chỉnh giá trị. Bảng~\ref{table:appendix_b1_archetypes} liệt kê các nguyên mẫu và Bảng~\ref{table:appendix_b1_rules} liệt kê các luật TraCI.

\begin{table}[H]
\centering
\small\setlength{\tabcolsep}{3pt}
\begin{tabular}{|p{3.2cm}|p{2.2cm}|p{1.4cm}|p{6.2cm}|}
\hline
\textbf{Nguyên mẫu} & \textbf{vClass cơ sở} & \textbf{Tỷ trọng} & \textbf{Tham số và nguồn} \\ \hline
\texttt{car\_\allowbreak saval} & \texttt{passenger} & 1,0 & \texttt{accel} 1,2; \texttt{decel} 4,8; \texttt{tau} 2,5; \texttt{minGap} 1,1. Cột ô tô Bảng 11 của Saval và cộng sự. Ô tô luôn dừng khi đèn đỏ nên không có luật vượt đỏ và các thuộc tính \texttt{jm*} giữ mặc định của SUMO \\ \hline
\texttt{moto\_\allowbreak saval} & \texttt{motorcycle} & 1,0 & \texttt{accel} 1,1; \texttt{decel} 4,6; \texttt{tau} 0,8; \texttt{minGap} 1,9; \texttt{minGapLat} 0,3. Cột xe máy Bảng 11. Khoảng cách ngang theo bài báo khoảng 0,1~m (làn xe máy 0,8~m gồm bề rộng 0,7~m và khoảng hở 0,1~m), nhỏ hơn cận dưới 0,3~m nên được kẹp về 0,3~m \\ \hline
\texttt{moto\_\allowbreak violator\_\allowbreak saval} & \texttt{custom1} & 1,0 & \texttt{accel} 1,1; \texttt{decel} 4,6; \texttt{tau} 0,8; \texttt{minGap} 1,9; \texttt{minGapLat} 0,3; \texttt{jmTimegapMinor} 1,0. Cùng giá trị với xe máy; \texttt{jmTimegapMinor} bằng mặc định của SUMO, là giả định vì bài báo không mô hình hóa rẽ trái trái phép \\ \hline
\end{tabular}
\caption{Các nguyên mẫu tài xế của B1 (đơn vị: m/s\textsuperscript{2} cho \texttt{accel}, \texttt{decel}; s cho \texttt{tau}, \texttt{jmTimegapMinor}; m cho \texttt{minGap}, \texttt{minGapLat})}
\label{table:appendix_b1_archetypes}
\end{table}

\begin{table}[H]
\centering
\small
\begin{tabular}{|p{3.2cm}|p{3.4cm}|p{3.0cm}|p{3.4cm}|}
\hline
\textbf{Luật} & \textbf{Áp dụng cho} & \textbf{Tham số} & \textbf{Nguồn} \\ \hline
\texttt{swarm\_\allowbreak red\_\allowbreak running} & \texttt{moto\_\allowbreak saval}, \texttt{moto\_\allowbreak violator\_\allowbreak saval} & $p_0 = 0{,}5$; $k = 0{,}0$; $p_{\max} = 0{,}5$; $\Delta t = 10{,}0$~s & Mục 4 của bài báo: xe máy ở hàng đầu chỉ dừng khi đèn đỏ với xác suất 0,5; $k = 0$ vì bài báo không có hiệu ứng đoàn; $\Delta t$ là giả định và không có tác dụng khi $k = 0$ \\ \hline
\texttt{early\_\allowbreak start} & \texttt{moto\_\allowbreak saval}, \texttt{moto\_\allowbreak violator\_\allowbreak saval} & $t_{\text{before\_green}} = 5{,}0$~s; $p = 0{,}0$; $v_{\max} = 2{,}0$~m/s & Không có trong bài báo (giả định): luật bị tắt ($p = 0$), chỉ giữ trong cấu trúc để các phương án tối ưu điều chỉnh \\ \hline
\end{tabular}
\caption{Các luật TraCI của B1}
\label{table:appendix_b1_rules}
\end{table}

B1 không có luật \texttt{illegal\_\allowbreak left\_\allowbreak box}, vì mô hình của Saval và cộng sự không có hành vi chờ giữa nút; vùng chờ của B1 vận hành theo cơ chế sẵn có của SUMO như B0. B1 cũng không có luật \texttt{wrong\_\allowbreak lane\_\allowbreak dodge} và \texttt{slow\_\allowbreak lane\_\allowbreak center}, và không có luật nào áp dụng cho ô tô.

\end{document}
